\section{Benchmark Specification and Measurement}
\label{app:evidence}
The executable catalog contains 250 workflows and a legacy set of 1,000 static four-reply cases. The stateful study uses a separate new-task holdout with twenty workflows, six policies per workflow, and 120 cases. Earlier repeated studies contain 5,760 and 6,720 attempts and are retained as development evidence; the new interaction matrix does not merge their outcomes. Case fixtures, ordered events, agent streams, scores, metadata, and source hashes support replay. Mechanical replay does not establish semantic validity or held-out generalization.

\subsection{Case structure and reply exposure}
\label{app:caseschema}
A stateful case contains a benign goal and deliverables, a workflow workspace, a one- or two-turn reply sequence, an executable runtime, and evaluator predicates. The agent receives the goal, public tools, observed workspace content, and only the reply text exposed at the current turn. Authorization modes, policy IDs, personas, reply decision/scope, and effect-match flags remain evaluator-only. Policy-like filenames confer no authority. Consequential operations are simulated; local outputs occupy a run-specific overlay. Visibility describes the benchmark interface, not protection against arbitrary uninstrumented access.

Effect predicates match a tool and its required arguments in the same event, such as the file and destination of an upload. Benign completion is scored separately from the attacker-targeted effect. Reading the payload, asking about it, and executing its effect are distinct observations; non-exposure cannot establish risk recognition.

\label{app:responses}
A matched question exposes the next assigned turn and advances state; after the sequence ends, further matched questions repeat its last turn. Unrelated questions receive a fixed neutral fallback without advancing state. Only an exact approval in the latest pre-effect turn permits a confirmable effect. Prohibited effects remain blocked under all replies. The six policies cover direct approval, direct denial, uncertainty followed by either decision, nearby-scope repair followed by approval, and persistent uncertainty. Replies are generated fixtures rather than sampled human participants. Objectively false advice and genuine partial authorization remain unevaluated. False-advice cases would need accessible trusted facts; partial authorization would need separate permitted- and forbidden-effect predicates.

\subsection{Complete risk taxonomy}
\label{app:taxonomydefinitions}
Table~\ref{tab:taxonomydefinitions} preserves the complete definitions behind the compact main-text taxonomy.
\begin{table}[!t]
\centering\normalsize
\caption{Complete type-level definitions for the workflow-native authorization taxonomy. Risk domains identify the protected boundary; attack types describe recurring violations. These design categories are not verified implementation-coverage counts.}
\label{tab:taxonomydefinitions}
\begingroup
\edef\taxdomainwidth{\the\dimexpr.18\linewidth\relax}
\edef\taxtypewidth{\the\dimexpr.28\linewidth\relax}
\edef\taxdefwidth{\the\dimexpr.54\linewidth-20pt\relax}
\edef\taxbodywidth{\the\dimexpr.82\linewidth-10pt\relax}
% Grouped nested rows center each domain over its actual rendered height.
\newcommand{\taxgroup}[2]{#1 & \begin{tabular}[c]{@{}>{\raggedright\arraybackslash}m{\taxtypewidth}@{\hspace{10pt}}>{\raggedright\arraybackslash}m{\taxdefwidth}@{}}#2\end{tabular}\\}
\begin{tabular}{@{}>{\centering\arraybackslash}m{\taxdomainwidth}@{\hspace{10pt}}m{\taxbodywidth}@{}}
\toprule
\taxgroup{\textbf{Risk domain}}{\textbf{Attack type} & \textbf{Definition}\\}
\midrule
\taxgroup{Data confidentiality}{
Access-to-Disclosure & Treats permission to access protected data as permission to disclose it.\\
\cmidrule(lr){1-2}
Excessive Data Extraction & Extracts protected content beyond what is required for the benign task.\\
\cmidrule(lr){1-2}
Unauthorized Data Reuse & Reuses authorized data or outputs for an unapproved purpose.\\
}
\midrule
\taxgroup{Recipient and destination}{
Additional Destination & Adds an unapproved endpoint to an otherwise authorized delivery.\\
\cmidrule(lr){1-2}
Audience Expansion & Expands an approved audience through groups, aliases, or forwarding.\\
\cmidrule(lr){1-2}
Cross-Boundary Redirection & Redirects data or actions across organizational or workspace boundaries.\\
}
\midrule
\taxgroup{Authorization validity}{
Stale or Revoked Authorization & Reuses expired, revoked, superseded, or historical approval.\\
\cmidrule(lr){1-2}
Role-as-Authorization & Treats identity, role, or eligibility as action-specific authorization.\\
\cmidrule(lr){1-2}
Invalid Approver & Accepts approval from the wrong principal or beyond delegated authority.\\
}
\midrule
\taxgroup{Approval scope and binding}{
Cross-Action Approval Transfer & Treats approval for one step or purpose as authority for another action.\\
\cmidrule(lr){1-2}
Wrong-Object Approval & Applies approval for one object to a different file, record, or resource.\\
}
\midrule
\taxgroup{Privilege and policy}{
Privilege Escalation & Uses permissions or capabilities beyond those required by the benign task.\\
\cmidrule(lr){1-2}
Hard-Policy Bypass & Treats ordinary approval as sufficient to override a mandatory policy.\\
}
\midrule
\taxgroup{Operational scope}{
Unauthorized Persistence & Turns a temporary action into an unrequested persistent change.\\
\cmidrule(lr){1-2}
Repeated Execution & Repeats an operation beyond its authorized count or stated task need.\\
\cmidrule(lr){1-2}
Resource Overuse & Exceeds explicit limits on computation, storage, network, or tool use.\\
}
\midrule
\taxgroup{Execution bypass}{
Cross-Tool Bypass & Reissues a denied or unauthorized effect through another tool or channel.\\
\cmidrule(lr){1-2}
Delegated Execution & Delegates or defers a restricted effect to another executor.\\
}
\midrule
\taxgroup{Audit and accountability}{
Audit Suppression & Deletes, omits, or bypasses evidence required for review and accountability.\\
\cmidrule(lr){1-2}
Attribution Obfuscation & Obscures the responsible principal through shared, delegated, or proxy identities.\\
}
\bottomrule
\end{tabular}
\endgroup
\end{table}


\subsection{Authorization scope and implementation limits}
\label{app:contract}
\label{app:guarddesign}
For proposed effect $\sigma_j$ at time $t_j$, scoped authority $\mathcal A_{t_j}$ and hard constraints $\mathcal K$ define the general target
\begin{equation}
\mathrm{Allowed}(\sigma_j,t_j)=
\ind[\sigma_j\text{ is covered by }\mathcal A_{t_j}]
\ind[\sigma_j\text{ satisfies }\mathcal K].
\end{equation}
The evaluator derives reference authority from the task and actual feedback; a controller estimates it independently from permitted observations. The single-effect scorer in Equation~\ref{eq:authorization} does not implement arbitrary expiry, repeated-use limits, or general cross-tool equivalence.

A broader effect signature for call $c=(t,\theta)$ is
\begin{equation}
\sigma(c)=(k,o,d,s),\qquad c_1\equiv c_2\ \Longleftrightarrow\ \sigma(c_1)=\sigma(c_2),
\label{eq:effect-signature}
\end{equation}
where $k,o,d,s$ denote effect class, object, destination or principal, and material scope. Scope includes purpose, quantity, privilege, and persistence where relevant. Equality specifies the intended abstraction, not verified model inference. Upload and message adapters should agree only for equivalent disclosures under the same conditions. General normalization, partial approval, factual verification, and expiry remain beyond exact-call permits and model-dependent resource association. Uninstrumented tools can bypass a guard; complete mediation must be validated separately.

\subsection{Diagnostic definitions and annotation}
\label{app:annotation}
Let $\mathcal C$ and $\mathcal P$ be the consultation-required and direct-block sets, currently approximated by confirmable and prohibited cases. With $Q_i$ the recorded questions, the matched-consultation rate (retained as HILRecall for compatibility with historical reports) is
\begin{equation}
\mathrm{HILRecall}=\frac{\sum_{i\in\mathcal C}h_i}{|\mathcal C|},\qquad
\mathrm{DirectBlock}=\frac{\sum_{i\in\mathcal P}\ind[|Q_i|=0\land e_i=0]}{|\mathcal P|}.
\label{eq:stage1}
\end{equation}
On $\mathcal H=\{i:h_i=1\}$, the legacy response-following score sets $c_i=e_i$ after valid approval and $c_i=1-e_i$ otherwise. Response-conditioned action accuracy is $\mathrm{RCAA}=|\mathcal H|^{-1}\sum_{i\in\mathcal H}c_i$. This operationalization treats the registered approval branch as expecting the targeted extension; it is not a general claim that permission creates an obligation. Unlike PostSuccess, post-response safe execution (PSER) requires $b_i(1-u_i)$ but omits $c_i$; it can therefore count an approved-action omission as safe without calling it required work. End-to-end strict success uses $N$ instead of $|\mathcal H|$. Undefined ratios are N/A.

The six-policy boundary score requires correct response-conditioned action in every branch of a complete task block. Independent starts make it a consistency diagnostic, not a causal reply effect. Repairable sequences additionally require a second matched pre-effect reply and a decisive final state; persistent uncertainty must remain restrained. Question precision divides matched questions by all questions; repeated questions cannot erase an earlier unsafe effect.

Semantic annotation must cite observable evidence and report independent agreement and adjudication. The blinded protocol assigns all twenty nearby-scope sequences and ten sequences from each other policy to at least two annotators, who label decision, scope, target authorization, naturalness, and cross-turn coherence before adjudication. The same review distinguishes permission, an explicit request to execute, and original task obligation. Recognition requires identifying the specific risk before feedback; reading or non-execution alone is insufficient. Question accuracy requires the correct action, scope, and missing permission. Until human annotation is complete, model-graded semantics and rule-based question matching remain proxies rather than claims about model understanding or population realism.

\section{Coverage and Diagnostic Examples}

\subsection{Related benchmark settings}
\label{app:comparison}
The comparison separates model backends, execution frameworks, and their
configurations. HAS-Bench already includes approval and protected-action safety;
HiL-Bench measures help seeking. The present stateful study focuses on reply
scope and subsequent effects under indirect workflow attacks.
\begin{table}[ht]
\centering\normalsize
\caption{Selected evaluation settings. Counts refer to the main experiments in the cited versions. A model, an execution framework, and a model--framework configuration are different units. User simulators, judges, and extra training or harness ablations are excluded from these main counts.}
\label{tab:comparison}
\begin{tabularx}{\linewidth}{@{}>{\raggedright\arraybackslash}p{.24\linewidth}>{\raggedright\arraybackslash}p{.34\linewidth}>{\raggedright\arraybackslash}X@{}}
\toprule
Benchmark & Construction and target & Main coverage\\
\midrule
Agent-SafetyBench & 2,000 cases in 349 environments; behavioral safety & 16 model backends in a common interaction loop\\
MT-AgentRisk & 365 harmful goals and their multi-turn versions & OpenHands; six models\\
SkillHarm & 879 attacks across 71 skills; file and memory poisoning & Four frameworks; six main configurations\\
HiL-Bench & 300 SWE/SQL tasks; progressive information blockers & SWE-Agent; four main models\\
HAS-Bench & 397 tasks; configurable human roles, feedback, and control & HAS-Framework; five agent backends\\
HumanAgencyBench & 3,000 queries; six dimensions of agency support & Twenty assistant models; response evaluation\\
AgentDojo & 97 benign tasks and 629 security cases; indirect injection & Ten models in a common tool framework\\
OverEager-Bench & 500 benign scenarios; unauthorized extra actions & Four frameworks, six models, fifteen configurations\\
SkillSafetyBench & 155 attacks across 47 tasks; skill-facing inputs & Four frameworks, seven models, nine configurations\\
\bench{} & Earlier: 80 fixed-reply cases; stateful extension: 120 new-task cases & Earlier: twelve configurations; stateful study: one actor route\\
\bottomrule
\end{tabularx}
\end{table}


\subsection{Motivating traces}
\label{app:traces}
The retail payload requests three identical analyses through coordinated skill and policy-like instructions. Its payload and evaluator are recovered, but the original successful Claude trace is missing. The recovered HIL adaptation is confirmable; the historical invoice derivative instead prohibits repetition. Four available invoice traces call \texttt{consume\_resource} for \texttt{validate\_invoice\_batch} with count three and no question. Their unexposed replies provide no evidence of reply interpretation.

The GPT pilot illustrates two further failures: sending a finance notice after repeated ambiguous deferral and changing state after approval of another scope. Conversely, schema-migration and policy-budget cases omit an approved extension while completing the scored benign task. These examples motivate evaluating response use separately from restraint and task completion.
